\chapter{Why attention: building a transformer layer}
\label{ch:attention}

\begin{objectives}
\begin{itemize}
\item Explain why chess punishes architectures that assume nearby squares matter most.
\item Build attention in six steps, from "average your neighbours" to the real formula.
\item Compute a complete attention operation by hand over three squares and check it.
\item Say what the query, key and value projections each mean, and why three are needed.
\item Explain the $\sqrt{d_k}$ with a numerical demonstration of what goes wrong without it.
\item Describe exactly what \texttt{BT3-768x15x24h} names, and count the attention matrices in
  one forward pass.
\end{itemize}
\end{objectives}

\section{The board as 64 tokens}

The network must turn a position into an evaluation. First it must represent the position as
numbers. Leela's transformer nets do this the way a language model represents a sentence:

\begin{keyidea}
\textbf{The board is a sequence of 64 tokens, one per square.} Each square is represented by
a vector of $d = 768$ numbers. Initially those numbers encode what is standing there (piece
type, colour) plus global facts (side to move, castling rights, repetition count, and several
previous positions --- the history of Chapter~\ref{ch:byhand} §7.6). Thereafter, each layer
\emph{rewrites} all 64 vectors, and what they encode is up to the training process.
\end{keyidea}

The whole position is therefore a $64 \times 768$ block of numbers, and the network's job is
to refine that block fifteen times and then read two answers out of it.

\begin{notationbox}
$\vv{x}_i \in \R^{768}$ is the vector at square $i$. $X \in \R^{64\times768}$ stacks all 64.
Superscripts denote layers: $X^{(0)}$ is the input encoding, $X^{(15)}$ the final one.
\end{notationbox}

\section{The problem with the obvious architecture}

Before transformers, chess networks were \emph{convolutional}: each square's new value is
computed from a small fixed neighbourhood --- itself and the eight squares around it --- and
the same rule is applied at every square. Stack enough of these and information spreads
across the board a step at a time.

That works, and Leela's earlier networks were built this way. But it is a poor match for
chess, for two reasons.

\begin{enumerate}
\item \textbf{Chess relations are long-range.} A bishop on c4 and the square f7 are related in
  the strongest possible way, and they are four squares apart with two squares in between
  which have nothing to do with it. A convolution must pass that relationship through several
  layers of neighbours, degrading it, when the relationship is direct.
\item \textbf{Chess relations are content-dependent.} Whether c4 and f7 are related \emph{at
  all} depends on what is on d5 and e6. A convolution's neighbourhood is fixed before the
  position is seen; the set of squares that matter is not.
\end{enumerate}

What we want instead is an operation where \textbf{each square decides for itself which other
squares to read from, based on what is on them}. That is attention.

\section{Building attention in six steps}

\begin{attempt}{1}{average everything}
Give every square the average of all 64 squares:
$\vv{x}_i' = \frac{1}{64}\sum_j \vv{x}_j$.
Information now travels any distance in one step --- the long-range problem is solved. But
every square receives the identical mush, and nothing depends on what is where. Useless, but
it establishes the shape: \textbf{a new value is a weighted sum of the others.}
\end{attempt}

\begin{attempt}{2}{fixed weights}
$\vv{x}_i' = \sum_j \alpha_{ij}\,\vv{x}_j$ with $\alpha_{ij}$ learned constants --- so the
network can learn that c4 should read heavily from f7. Better, and it can encode geometry.
But $\alpha$ is fixed once training ends: c4 reads from f7 with the same intensity whether the
diagonal is open, blocked, or the bishop is not there at all. \textbf{The weights must depend
on the content of the squares, not just their names.}
\end{attempt}

\begin{attempt}{3}{weights from similarity}
Let the weight between two squares be their \emph{similarity}, measured by the dot product:
\[
  e_{ij} \;=\; \vv{x}_i \cdot \vv{x}_j \;=\; \sum_{c=1}^{d} x_{ic}\,x_{jc}.
\]
Recall what a dot product does: it is large and positive when two vectors point the same way,
zero when they are perpendicular, negative when opposed. So $e_{ij}$ is large when squares $i$
and $j$ carry "compatible" content --- and what counts as compatible is learned.

\begin{worked}{dot products as a compatibility test}
Three two-dimensional vectors: $\vv{a} = (1, 0)$, $\vv{b} = (0.9, 0.4)$,
$\vv{c} = (-0.7, 0.7)$.
\[
\vv{a}\cdot\vv{b} = 0.9,\qquad \vv{a}\cdot\vv{c} = -0.7,\qquad \vv{b}\cdot\vv{c} = -0.35 .
\]
$\vv{a}$ and $\vv{b}$ point nearly the same way: strongly compatible. $\vv{a}$ and $\vv{c}$
point away from each other: incompatible. \emph{(Illustrative.)}
\end{worked}
\end{attempt}

\begin{attempt}{4}{make the weights a distribution --- softmax}
The scores $e_{ij}$ can be negative and do not sum to anything in particular, so they are not
usable as mixing weights. Convert them into a probability distribution over the 64 squares
with the \textbf{softmax}:
\begin{equation}
  \label{eq:softmax}
  \alpha_{ij} \;=\; \frac{e^{\,e_{ij}}}{\sum_{k=1}^{64} e^{\,e_{ik}}} .
\end{equation}
Exponentiating makes everything positive; dividing by the sum makes the row add to 1. The
result is literally \emph{how square $i$ divides its attention over the board}.

\begin{worked}{softmax by hand}
Scores $(2.0,\ 1.0,\ 0.0,\ -1.0)$. Exponentiate: $7.389,\ 2.718,\ 1.000,\ 0.368$. Sum:
$11.475$. Divide:
\[
  (0.644,\ 0.237,\ 0.087,\ 0.032) .
\]
Sums to 1.000. \checkmark\ Note the shape: a score gap of 1.0 became a ratio of about
$2.7\times$, and the lowest score still gets $3\%$ --- softmax is soft. That is deliberate:
a hard maximum would have no gradient to learn from.
\end{worked}
\end{attempt}

\begin{attempt}{5}{separate the three roles --- query, key, value}
Using the raw $\vv{x}_i$ for everything conflates three different questions. For a rook on an
open file:
\begin{itemize}
\item \textbf{what am I looking for?} squares along my file, and whatever obstructs them;
\item \textbf{what do I advertise to others?} "there is a rook here, on this file";
\item \textbf{what do I pass on when read?} perhaps something else again --- the fact that
  this file is contested.
\end{itemize}
These are three distinct pieces of information and one vector cannot be all of them at once.
So project the vector three ways, with three learned matrices:
\[
  \vv{q}_i = W_Q\,\vv{x}_i \quad(\text{query}),\qquad
  \vv{k}_i = W_K\,\vv{x}_i \quad(\text{key}),\qquad
  \vv{v}_i = W_V\,\vv{x}_i \quad(\text{value}).
\]
The score becomes $e_{ij} = \vv{q}_i\cdot\vv{k}_j$ --- \emph{does what $i$ is looking for
match what $j$ advertises?} --- and the thing that is mixed is $\vv{v}_j$.

Notice this makes attention \emph{asymmetric}: $i$ may attend strongly to $j$ while $j$
ignores $i$. Chess relations are asymmetric (a pinned knight and the pinning bishop are not in
symmetric situations), so this is a feature.
\end{attempt}

\begin{attempt}{6}{scale the scores}
One numerical problem remains. The dot product of two $d$-dimensional vectors is a sum of $d$
terms, so its typical size grows like $\sqrt{d}$. With $d = 768$ the scores become large, and
softmax of large numbers is brutal:

\begin{worked}{what large scores do to softmax}
Two scores differing by 2.0:
\begin{center}
\begin{tabular}{@{}lll@{}}
\toprule
Scores & softmax & \\
\midrule
$(2.0,\ 0.0)$   & $(0.881,\ 0.119)$ & a preference \\
$(20,\ 18)$     & $(0.881,\ 0.119)$ & the same --- only the difference matters \\
$(200,\ 180)$   & $(1.000,\ 2\times10^{-9})$ & effectively one-hot \\
\bottomrule
\end{tabular}
\end{center}
The third row is the danger: attention collapses onto a single square and the gradient
through the other entries vanishes, so the network cannot learn to change its mind. Dividing
the scores by $\sqrt{d_k}$ keeps them in the first regime regardless of dimension.
\end{worked}
\end{attempt}

\begin{finalform}[Scaled dot-product attention]
\begin{equation}
  \label{eq:attention}
  \boxed{\;
  \text{Attention}(Q,K,V) \;=\; \softmax\!\left(\frac{QK^{\top}}{\sqrt{d_k}}\right)V
  \;}
\end{equation}
where $Q = XW_Q$, $K = XW_K$, $V = XW_V$ are $64\times d_k$, and $QK^{\top}$ is the
$64\times64$ matrix of scores. Row $i$ of $\softmax(QK^\top/\sqrt{d_k})$ is \textbf{square
$i$'s attention distribution over the board}, and it is the object rendered as a heatmap in
your frontend.
\end{finalform}

\section{A complete attention computation, by hand}

\begin{worked}{three squares, two dimensions}
Three tokens (think: three squares of interest), $d = d_k = 2$.
\emph{(Illustrative numbers chosen for clean arithmetic.)}
\[
\vv{x}_1 = (1,\,0), \qquad \vv{x}_2 = (0,\,1), \qquad \vv{x}_3 = (1,\,1).
\]
Take
\[
W_Q = W_K = \begin{pmatrix}1&0\\0&1\end{pmatrix},
\qquad
W_V = \begin{pmatrix}1&0\\0&2\end{pmatrix}.
\]
Then $\vv{q}_i = \vv{k}_i = \vv{x}_i$, and
\[
\vv{v}_1 = (1,0),\qquad \vv{v}_2 = (0,2),\qquad \vv{v}_3 = (1,2).
\]

\textbf{Step 1 --- raw scores for square 1} ($\vv{q}_1 = (1,0)$):
\[
e_{11} = (1,0)\cdot(1,0) = 1,\qquad
e_{12} = (1,0)\cdot(0,1) = 0,\qquad
e_{13} = (1,0)\cdot(1,1) = 1 .
\]

\textbf{Step 2 --- scale} by $\sqrt{d_k} = \sqrt{2} = 1.4142$:
\[
(0.7071,\quad 0,\quad 0.7071).
\]

\textbf{Step 3 --- softmax:} $e^{0.7071} = 2.0281$, $e^{0} = 1$, $e^{0.7071} = 2.0281$;
sum $= 5.0562$:
\[
\boldsymbol{\alpha}_1 = (0.4011,\ 0.1978,\ 0.4011).
\]
Square 1 gives $40\%$ of its attention to itself, $40\%$ to square 3, and $20\%$ to square 2
--- because squares 1 and 3 share a component and square 2 does not.

\textbf{Step 4 --- mix the values:}
\[
\vv{x}_1' = 0.4011(1,0) + 0.1978(0,2) + 0.4011(1,2)
         = (0.4011 + 0.4011,\ \ 0.3956 + 0.8022)
         = (0.8022,\ 1.1978).
\]

Square 1 entered the layer as $(1,0)$ and leaves carrying $(0.80,\,1.20)$ --- it has absorbed
information from the squares it found relevant. That is the whole operation.

\smallskip
\emph{Check your understanding:} compute $\vv{x}_2'$ yourself (Exercise~\ref{ex:ch11:byhand}).
You should find square 2 attends most to square 3.
\end{worked}

\section{From one attention to a transformer layer}

Real layers wrap the operation in three pieces of standard machinery.

\paragraph{Multi-head.} One attention pattern per layer is not enough --- a square may need to
track its attackers \emph{and} its defenders \emph{and} the pawn structure simultaneously.
So run $H$ independent attentions in parallel, each with its own $W_Q, W_K, W_V$ of reduced
width, and concatenate the results:
\begin{equation}
  \label{eq:multihead}
  \text{MultiHead}(X) = \bigl[\text{head}_1 \,\|\, \text{head}_2 \,\|\, \cdots \,\|\,
  \text{head}_H\bigr]\,W_O .
\end{equation}
In \texttt{BT3-768x15x24h}, $H = 24$. Each head produces its own $64\times64$ attention
matrix, and the heads specialise --- which is what makes Chapter~\ref{ch:attmaps}'s warning
about averaging them so important.

\paragraph{Residual connections.} Each sub-layer adds to its input rather than replacing it:
$X \leftarrow X + \text{sublayer}(X)$. This means a layer only has to compute a
\emph{correction}, and it gives gradients a short path back through fifteen layers. Without
residuals, deep stacks are close to untrainable.

\paragraph{Layer normalisation and the feed-forward block.} After attention, each square's
vector is renormalised (subtract the mean, divide by the standard deviation, then apply a
learned scale and shift) to keep numbers in a sane range. Then each square independently
passes through a two-layer network of the Chapter~\ref{ch:learned} kind --- the
\emph{feed-forward} block, typically expanding to $4d$ and back.

\begin{finalform}[One encoder layer]
\begin{align}
  Z &= \LN\bigl(X + \text{MultiHead}(X)\bigr), \label{eq:enc1}\\
  X' &= \LN\bigl(Z + \text{FFN}(Z)\bigr), \qquad
  \text{FFN}(\vv{z}) = W_2\,\relu(W_1\vv{z} + \vv{b}_1) + \vv{b}_2 . \label{eq:enc2}
\end{align}
\end{finalform}

The division of labour is clean and worth remembering:

\begin{keyidea}
\textbf{Attention moves information between squares. The feed-forward block thinks about one
square at a time.} Every layer does both: gather what is relevant from elsewhere on the board,
then process it locally. Fifteen rounds of that is what \texttt{BT3} does to a position before
it says anything.
\end{keyidea}

\section{Reading the name of the network}

\begin{center}
\texttt{BT3-768x15x24h}
\end{center}

\begin{center}
\begin{tabular}{@{}lll@{}}
\toprule
\texttt{768} & embedding width $d$ & each square is 768 numbers \\
\texttt{15}  & layers & fifteen rounds of Equations \eqref{eq:enc1}--\eqref{eq:enc2} \\
\texttt{24h} & heads & 24 attention patterns per layer \\
\bottomrule
\end{tabular}
\end{center}

\begin{worked}{how many attention matrices in one evaluation?}
Each head at each layer produces one $64\times64$ matrix:
\[
  15 \times 24 = 360 \text{ matrices}, \qquad 360 \times 64 \times 64 = 1{,}474{,}560
  \text{ attention weights}
\]
for a single position. Your saliency display averages some of these into one $8\times8$
picture. Chapter~\ref{ch:attmaps} is about what that averaging destroys --- but note the
compression ratio now: 1.47 million numbers into 64.
\end{worked}

And a scale note, so the cost of the interpretability work in Part~\textsc{v} is not a
surprise: one forward pass of BT3 on this machine's CPU takes on the order of a second, which
is why every real number in this book came from the smaller \texttt{791556} net. The
difference is not conceptual --- same architecture family, fewer parameters.

\begin{repobox}[title={in this repository}]
\texttt{backend/neural\_vision.py} attaches forward hooks to these modules through
\texttt{lczerolens}, which is how the attention heatmap is produced: it reaches inside
Equation~\eqref{eq:multihead} and captures $\softmax(QK^\top/\sqrt{d_k})$ before it is
multiplied by $V$. The memory note in your project about black-to-move orientation is a
statement about which square index $i$ corresponds to which board square --- a detail that
Equation~\eqref{eq:attention} is entirely silent about, and which therefore has to be got
right by hand.
\end{repobox}

\begin{exercises}

\exhead[dot products]{ch11:dots} Compute $\vv{a}\cdot\vv{b}$ for $\vv{a} = (2,-1,0)$ and
$\vv{b} = (1,3,4)$; then for $\vv{b} = (2,-1,0)$. Which pair is "more compatible", and what
does a negative dot product mean in attention terms?

\exhead[softmax]{ch11:softmax} Compute the softmax of $(3,\ 1,\ 0,\ 0)$ and of
$(6,\ 2,\ 0,\ 0)$. Which is sharper? Now compute the softmax of $(1003,\ 1001,\ 1000,\ 1000)$
and comment on what the identical answer to the first tells you about softmax.

\exhead[attention by hand]{ch11:byhand} Using the setup of §11.4, compute $\vv{x}_2'$ and
$\vv{x}_3'$ in full: raw scores, scaled scores, softmax, mixed value. Which square does
square 2 attend to most, and does that make sense from the vectors?

\exhead[asymmetry]{ch11:asym} With $W_Q = W_K$ as in §11.4, is $\alpha_{ij} = \alpha_{ji}$?
Compute $\alpha_{13}$ and $\alpha_{31}$ to check. Now argue why real networks make
$W_Q \neq W_K$, using a chess relation that is genuinely one-directional.

\exhead[the scaling]{ch11:scaling} Suppose scores are typically of size $\sqrt{d_k}$. Compute
the softmax of $(\sqrt{768},\ 0)$ without scaling, and of $(\sqrt{768}/\sqrt{768},\ 0)$ with
it. State in one sentence what training would look like in the first case.

\exhead[why three matrices]{ch11:threemat} Give a chess example, other than the rook one in
§11.3, where "what I am looking for" and "what I advertise" are clearly different. Then say
what would break if $W_Q$, $W_K$ and $W_V$ were forced to be equal.

\exhead[convolution vs attention]{ch11:conv} A convolutional layer mixes each square with its
8 neighbours. How many layers are needed before a1 can be influenced by h8? How many
attention layers? What does the answer suggest about which architecture better suits a
bishop?

\exhead[counting]{ch11:counting} How many attention weights does one BT3 forward pass produce
(§11.5)? If you displayed one $8\times8$ heatmap per head per layer, how many pictures would
that be, and how would you choose which to show a user?

\exhead[what a row means]{ch11:row} Row $i$ of the attention matrix sums to 1. What does
column $j$ sum to, and why is that not 1? Which of the two --- "where did e4 look?" or "who
looked at e4?" --- is a row, and which would you rather display over a chessboard?

\exhead[residuals]{ch11:residual} Explain in two sentences why $X + \text{sublayer}(X)$ makes
a 15-layer stack trainable, and what it implies about how much any one layer changes the
representation. Relate your answer to the idea of a layer computing a "correction".

\exhead[design]{ch11:design} Your tool wants to show "what the network is looking at" for a
single square the user clicks. Which row or column, of which head, of which layer, would you
display? List the choices you must make, and note which of them Chapter~\ref{ch:attmaps} will
tell you are dangerous.

\end{exercises}
